%%% LaTeX-Vorlage Version 1.8 %%%

% Grundlegende Dokumenteneigenschaften gemäß DHBW-Vorgaben
\documentclass[a4paper,fontsize=11pt,oneside,parskip=half,headings=normal]{scrreprt} 
% \usepackage{showframe} % nur für Kontrolle der Ränder 

%%% Präambel einbinden (mit Festlegungen gemäß DHBW-Vorgaben) %%%
%%% Präambel %%%
% hier sollten keine Änderungen erforderlich sein
%
\usepackage[utf8]{inputenc}   % Zeichencodierung UTF-8 für Eingabe-Dateien
\usepackage[T1]{fontenc}      % Darstellung von Umlauten im PDF

\usepackage{listings}         % für Einbindung von Code-Listings
\lstset{numbers=left,numberstyle=\tiny,numbersep=5pt,texcl=true}
\lstset{literate=             % erlaubt Sonderzeichen in Code-Listings 
{Ö}{{\"O}}1
{Ä}{{\"A}}1
{Ü}{{\"U}}1
{ß}{{\ss}}2
{ü}{{\"u}}1
{ä}{{\"a}}1
{ö}{{\"o}}1
{€}{{\euro}}1
}

\usepackage[
  inner=35mm,outer=15mm,top=25mm,
  bottom=20mm,foot=12mm,includefoot
]{geometry}                 % Einstellungen für Ränder

\usepackage[ngerman]{babel} % Spracheinstellungen Deutsch
\usepackage[babel,german=quotes]{csquotes} % deutsche Anf.zeichen
\usepackage{enumerate}      % anpassbare Nummerier./Aufz.
\usepackage{graphicx}       % Einbinden von Grafiken
\usepackage[onehalfspacing]{setspace} % anderthalbzeilig

\usepackage{blindtext}      % Textgenerierung für Testzwecke
\usepackage{color}          % Verwendung von Farbe 

\usepackage{acronym}        % für ein Abkürzungsverzeichnis

\usepackage[                % Biblatex
  backend=biber,
  bibstyle=_dhbw_authoryear,maxbibnames=99,
  citestyle=authoryear,     
  uniquename=true, useprefix=true,
  bibencoding=utf8]{biblatex}
%kein Punkt am Ende bei \footcite
%http://www.golatex.de/footcite-ohne-punkt-am-schluss-t4865.html
\renewcommand{\bibfootnotewrapper}[1]{\bibsentence#1}


%Reihenfolge der Autorennamen
%   
% http://golatex.de/viewtopic,p,80448.html#80448
% Argumente: siehe http://texwelt.de/blog/modifizieren-eines-biblatex-stils/
\DeclareNameFormat{sortname}{% Bibliographie
  \ifnum\value{uniquename}=0 % Normalfall
    \ifuseprefix%
      {%
         \usebibmacro{name:family-given}
           {\namepartfamily}
           {\namepartgiveni}
           {\namepartprefix}
           {\namepartsuffixi}%
       }
      {%
         \usebibmacro{name:family-given}
           {\namepartfamily}
           {\namepartgiveni}
           {\namepartprefixi}
           {\namepartsuffixi}%
       }%
  \fi
  \ifnum\value{uniquename}=1% falls nicht eindeutig, abgek. Vorname 
      {%
         \usebibmacro{name:family-given}
           {\namepartfamily}
           {\namepartgiveni}
           {\namepartprefix}
           {\namepartsuffix}%
       }%
  \fi
  \ifnum\value{uniquename}=2% falls nicht eindeutig, ganzer Vorname 
      {%
         \usebibmacro{name:family-given}
           {\namepartfamily}
           {\namepartgiven}
           {\namepartprefix}
           {\namepartsuffix}%
       }%
  \fi   
  \usebibmacro{name:andothers}}

\DeclareNameFormat{labelname}{% für Zitate
  \ifnum\value{uniquename}=0 % Normalfall
    \ifuseprefix%
      {%
         \usebibmacro{name:family-given}
           {\namepartfamily}
           {\empty}
           {\namepartprefix}
           {\namepartsuffixi}%
       }
      {%
         \usebibmacro{name:family-given}
           {\namepartfamily}
           {\empty}
           {\namepartprefixi}
           {\namepartsuffixi}%
       }%
  \fi
  \ifnum\value{uniquename}=1% falls nicht eindeutig, abgek. Vorname 
      {%
         \usebibmacro{name:family-given}
           {\namepartfamily}
           {\namepartgiveni}
           {\namepartprefix}
           {\namepartsuffix}%
       }%
  \fi
  \ifnum\value{uniquename}=2% falls nicht eindeutig, ganzer Vorname 
      {%
         \usebibmacro{name:family-given}
           {\namepartfamily}
           {\namepartgiven}
           {\namepartprefix}
           {\namepartsuffix}%
       }%
  \fi   
  \usebibmacro{name:andothers}}
      
  
\DeclareFieldFormat{extrayear}{% = the 'a' in 'Jones 1995a'
  \iffieldnums{labelyear}
    {\mknumalph{#1}}
    {\mknumalph{#1}}}        

\renewcommand*{\multinamedelim}{\addslash}
\renewcommand*{\finalnamedelim}{\addslash}
\renewcommand*{\multilistdelim}{\addslash}
\renewcommand*{\finallistdelim}{\addslash}

\renewcommand{\nameyeardelim}{~}

% Literaturverzeichnis: Doppelpunkt zwischen Name (Jahr): Rest 
% http://de.comp.text.tex.narkive.com/Tn1HUIXB/biblatex-authoryear-und-doppelpunkt
\renewcommand{\labelnamepunct}{\addcolon\addspace}

% damit die Darstellung für Vollzitate von Primärquellen in 
% Fußnoten später auf "nicht fett" geändert werden kann 
% (nur für Zitate von Sekundärliteratur relevant)
\newcommand{\textfett}[1]{\textbf{#1}}

% für Zitate von Sekundärliteratur:
\newcommand{\footcitePrimaerSekundaer}[4]{%
  \renewcommand{\textfett}[1]{##1}%
  \footnote{\fullcite[#2]{#1}, zitiert nach \cite[#4]{#3}}%  
  \renewcommand{\textfett}[1]{\textbf{##1}}%
}

% Im Literaturverzeichnis: Autor (Jahr) fett
\renewbibmacro*{author}{%
  \ifboolexpr{%
    test \ifuseauthor%
    and
    not test {\ifnameundef{author}}
  }
    {\usebibmacro{bbx:dashcheck}
       {\bibnamedash}
       {\usebibmacro{bbx:savehash}%
        \textfett{\printnames{author}}%
        \iffieldundef{authortype}
          {\setunit{\addspace}}
          {\setunit{\addcomma\space}}}%
     \iffieldundef{authortype}
       {}
       {\usebibmacro{authorstrg}%
        \setunit{\addspace}}}%
    {\global\undef\bbx@lasthash
     \usebibmacro{labeltitle}%
     \setunit*{\addspace}}%
  \textfett{\usebibmacro{date+extrayear}}}

% Sonderfall: Quelle ohne Autor, aber mit Herausgeber
% Name des Herausgebers wird fett gedruckt
\renewbibmacro*{bbx:editor}[1]{%
  \ifboolexpr{%
    test \ifuseeditor%
    and
    not test {\ifnameundef{editor}}
  }
    {\usebibmacro{bbx:dashcheck}
       {\bibnamedash}
       {\textfett{\printnames{editor}}%
        \setunit{\addcomma\space}%
        \usebibmacro{bbx:savehash}}%
     \usebibmacro{#1}%
     \clearname{editor}%
     \setunit{\addspace}}%
    {\global\undef\bbx@lasthash
     \usebibmacro{labeltitle}%
     \setunit*{\addspace}}%
  \textfett{\usebibmacro{date+extrayear}}}

% Anpassungen für deutsche Sprache
\DefineBibliographyStrings{ngerman}{%
	nodate = {{o.J.}},
	urlseen = {{Abruf:}},
	ibidem = {{ebenda}}
}

% keine Anführungszeichen beim Titel im Literaturverzeichnis
\DeclareFieldFormat[article,book,inbook,inproceedings,manual,misc,phdthesis,thesis,online,report]{title}{#1\isdot}

\newcommand{\literaturverzeichnis}{%
% nur Literaturverzeichnis
% (als eigenes Kapitel)
\phantomsection
\addcontentsline{toc}{chapter}{Literaturverzeichnis}
\spezialkopfzeile{Literaturverzeichnis}
\defbibheading{lit}{\chapter*{Literaturverzeichnis}}
\label{chapter:quellen}
\printbibliography[heading=lit,notkeyword=ausblenden]
} % mit DHBW-spezifischen Einstellungen

\usepackage{hyperref}       % URL-Formatierung, klickbare Verweise

\usepackage{tocloft}        % für Verzeichnis der Anhänge

\usepackage{multirow}

\newcounter{anhcnt}
\setcounter{anhcnt}{0}
\newlistof{anhang}{app}{}

\newcommand{\anhang}[1]{%
  \refstepcounter{anhcnt}
  \setcounter{anhteilcnt}{0}
  \section*{Anhang \theanhcnt: #1}
  \addcontentsline{app}{section}{\protect\numberline{Anhang \theanhcnt}#1}\par
}

\newcounter{anhteilcnt}
\setcounter{anhteilcnt}{0}

\newcommand{\anhangteil}[1]{%
	\refstepcounter{anhteilcnt}
	\subsection*{Anhang~\arabic{anhcnt}/\arabic{anhteilcnt}: #1}
	\addcontentsline{app}{subsection}{\protect\numberline{Anhang \theanhcnt/\arabic{anhteilcnt}}#1}\par
}

\renewcommand{\theanhteilcnt}{Anhang \theanhcnt/\arabic{anhteilcnt}}

% vgl. S. 4 Paket-Beschreibung tocloft 	
% Einrückungen für Anhangverzeichnis
\makeatletter
\newcommand{\abstaendeanhangverzeichnis}{
\renewcommand*{\l@section}{\@dottedtocline{1}{0em}{5.5em}}
\renewcommand*{\l@subsection}{\@dottedtocline{2}{2.3em}{6.5em}}
}
\makeatother

% Abbildungs- und Tabellenverzeichnis
% Bezeichnungen
\renewcaptionname{ngerman}{\figurename}{Abb.}
\renewcaptionname{ngerman}{\tablename}{Tab.}
% Einrückungen
\makeatletter
\renewcommand*{\l@figure}{\@dottedtocline{1}{0em}{2.3em}}
\renewcommand*{\l@table}{\@dottedtocline{1}{0em}{2.3em}}
\makeatother


\usepackage{chngcntr}                % fortlaufende Zähler für Fußnoten, Abbildungen und Tabellen
\counterwithout{figure}{chapter}
\counterwithout{table}{chapter}
\counterwithout{footnote}{chapter}

\usepackage[automark]{scrlayer-scrpage} 
%% Definitionen für Kopf- und Fußzeile auf normalen Seiten
\defpagestyle{kopfzeile}
{% Kopfdefinition
  (\textwidth,0pt)    % Länge der oberen Linie,Dicke der oberen Linie       
  {} % Definition für linke Seiten im doppelseitigen Layout
  {} % Definition für rechte Seiten im doppelseitigen Layout      
  {  % Definition für Seiten im einseitigen Layout
	\makebox[0pt][l]{\rightmark}% 
	\makebox[\linewidth]{}% 
  }        
  (\textwidth, 0.4pt) % Untere Linienlänge, Untere Liniendicke
}
{% Fußdefinition
  (\textwidth,0pt)    % Obere Linienlänge, Obere Liniendicke
  {} % Definition für linke Seiten im doppelseitigen Layout
  {} % Definition für rechte Seiten im doppelseitigen Layout
  {  % Definition für Seiten im einseitigen Layout
    \makebox[\linewidth]{}%
    \makebox[0pt][r]{\pagemark}%
  }
  (\textwidth, 0pt)   % Länge der unteren Linie,Dicke der unteren Linie
}

%% Definitionen für Kopf- und Fußzeile auf ersten Seiten eines Kapitels
\defpagestyle{kapitelkopfzeile}
{% Kopfdefinition
  (\textwidth,0pt)    % Länge der oberen Linie,Dicke der oberen Linie       
  {} % Definition für linke Seiten im doppelseitigen Layout
  {} % Definition für rechte Seiten im doppelseitigen Layout      
  {}  % Definition für Seiten im einseitigen Layout
  (\textwidth, 0pt) % Untere Linienlänge, Untere Liniendicke
}
{% Fußdefinition
  (\textwidth,0pt)    % Obere Linienlänge, Obere Liniendicke
  {} % Definition für linke Seiten im doppelseitigen Layout
  {} % Definition für rechte Seiten im doppelseitigen Layout
  {  % Definition für Seiten im einseitigen Layout
    \makebox[\linewidth]{}%
    \makebox[0pt][r]{\pagemark}%
  }
  (\textwidth, 0pt)   % Länge der unteren Linie,Dicke der unteren Linie
}

%% Definitionen für Kopf- und Fußzeile im Anhang und bei Quellenverzeichnisse
\newcommand{\spezialkopfzeileBezeichnung}{}
\defpagestyle{spezialkopfzeile}
{% Kopfdefinition
  (\textwidth,0pt)    % Länge der oberen Linie,Dicke der oberen Linie       
  {} % Definition für linke Seiten im doppelseitigen Layout
  {} % Definition für rechte Seiten im doppelseitigen Layout      
  {  % Definition für Seiten im einseitigen Layout
	\makebox[0pt][l]{\spezialkopfzeileBezeichnung}% 
	\makebox[\linewidth]{}% 
  }        
  (\textwidth, 0.4pt) % Untere Linienlänge, Untere Liniendicke
}
{% Fußdefinition
  (\textwidth,0pt)    % Obere Linienlänge, Obere Liniendicke
  {} % Definition für linke Seiten im doppelseitigen Layout
  {} % Definition für rechte Seiten im doppelseitigen Layout
  {  % Definition für Seiten im einseitigen Layout
    \makebox[\linewidth]{}%
    \makebox[0pt][r]{\pagemark}%
  }
  (\textwidth, 0pt)   % Länge der unteren Linie,Dicke der unteren Linie
}
            
\newcommand\spezialkopfzeile[1]{%
  \renewcommand\spezialkopfzeileBezeichnung{#1}
  \pagestyle{spezialkopfzeile}
}
                
% Standard-Pagestyle auswählen
\pagestyle{kopfzeile}

% keine Kopfzeile anzeigen auf Seiten, auf denen ein 
% Kapitel beginnt oder das Inhalts-/Abbildungs-/Tabellenverzeichnis steht 
\renewcommand{\chapterpagestyle}{kapitelkopfzeile}
\tocloftpagestyle{kapitelkopfzeile}

		 % für schöne Kopfzeilen 

\usepackage{textcomp}            % erlaubt EUR-Zeichen in Eingabedatei
\usepackage{eurosym}             % offizielles EUR-Symbol in Ausgabe
\renewcommand{\texteuro}{\euro}  % ACHTUNG: nach hyperref aufrufen!

\usepackage{scrhack}             % stellt Kompatibilität zw. KOMA-Script
                                 % (scrreprt) und anderen Paketen her
                                 
% Anpassung der Abstände bei Kapitelüberschriften
% (betrifft auch Inhalts-, Abbildungs- und Tabellenverzeichnis)
\renewcommand*\chapterheadstartvskip{\vspace*{-\topskip}}
\newcommand{\myBeforeTitleSkip}{1mm}
\newcommand{\myAfterTitleSkip}{10mm}
\setlength\cftbeforetoctitleskip{\myBeforeTitleSkip}
\setlength\cftbeforeloftitleskip{\myBeforeTitleSkip}
\setlength\cftbeforelottitleskip{\myBeforeTitleSkip}

\setlength\cftaftertoctitleskip{\myAfterTitleSkip}
\setlength\cftafterloftitleskip{\myAfterTitleSkip}
\setlength\cftafterlottitleskip{\myAfterTitleSkip}                                                            
%%% Ende der Präambel %%%


%%% Name der eigenen Literatur-Datenbank (ggf. anpassen) %%%
\bibliography{includes/literatur-datenbank.bib}

\begin{document}
%%% Deckblatt einbinden %%% 
% Anpassungen nötig (Name, Titel etc.)
% HIER EDITIEREN: 
% Typ der Arbeit (für Deckblatt und ehrenwörtliche Erklärung)
% - bitte Zutreffendes auswählen
\newcommand{\typMeinerArbeit}{1. Projektarbeit} 
%\newcommand{\typMeinerArbeit}{2. Projektarbeit} 
%\newcommand{\typMeinerArbeit}{Seminararbeit} 
% \newcommand{\typMeinerArbeit}{Bachelorarbeit} 

% Thema der Arbeit (für ehrenwörtliche Erklärung, ohne Umbrüche)
% HIER EDITIEREN: 
\newcommand{\themaMeinerArbeit}{Was ist ein Chatbot?}

% Vorname, Name der Autorin/des Autors (für Titelseite und Metadaten)
% HIER EDITIEREN:
\newcommand{\meinName}{Max Weiberle}

\thispagestyle{empty}

\begin{spacing}{1}
\begin{center}	
~\vspace{0mm}

% HIER EDITIEREN: Titel der Arbeit
{\sffamily
\LARGE  
% \Large  % bei sehr langen Titeln ggf. etwas kleinere Schriftart wählen
\textbf{Was ist ein Chatbot?}

\bigskip
\textbf{}
}


\vspace{15mm}

% Typ wird automatisch eingefügt (oben festlegen)
{\Large \typMeinerArbeit}

\vspace{1cm}

% HIER ggf. EDITIEREN
vorgelegt am 18. Mai 2022

\vspace{15mm}

Fakultät Wirtschaft
\medskip

Studiengang Wirtschaftsinformatik
\medskip

% HIER EDITIEREN: Kurs eintragen
Kurs WWI2021G

\vspace{10mm}

von

\vspace{10mm}

% Vorname und Name wird automatisch eingefügt (oben festlegen) 
{\large\textsc{\meinName}}

\vspace{10mm}
\end{center}

\vfill

% HIER EDITIEREN: Name des Unternehmens, Name der Betreuerin/des Betreuers
\begin{tabular}{ll}
Betreuer in der Ausbildungsstätte: & DHBW Stuttgart: \\
\hspace{0.4\linewidth} & \\
Komm.ONE & Prof., Dr., Manfred Sander \\
Stefan Fuchs \\
Studierendenbetreuer \\
\\
Unterschrift der Betreuerin/des Betreuers \\
\end{tabular}


\vspace{1cm}
%(etwas Platz für die Unterschrift der Betreuerin/des Betreuers aus der Ausbildungsstätte)
\end{spacing}

% falls ein Vertraulichkeitsvermerk erforderlich ist,
% die Kommentarzeichen in den nachfolgenden Zeilen entfernen:
 
%\begin{center}
%\small
%\textbf{Vertraulichkeitsvermerk}:
%Der Inhalt dieser Arbeit darf weder als Ganzes noch in Auszügen \\
%Personen außerhalb des Prüfungs- und Evaluationsverfahrens zugänglich gemacht werden, sofern keine anders lautende Genehmigung des Dualen Partners vorliegt. 
%\end{center}

% Meta-Daten für PDF-Datei basierend auf obigen Angaben
\hypersetup{pdftitle={\themaMeinerArbeit}}
\hypersetup{pdfauthor={\meinName}}
\hypersetup{pdfsubject={\typMeinerArbeit\ DHBW Stuttgart \the\year}}


%%% Umstellung der Seiten-Nummerierung auf i, ii, iii ... %%%
\pagenumbering{Roman} 

%%% Abstract einbinden (optionale Kurzfassung Ihrer Arbeit) %%%
\begin{abstract}
\thispagestyle{kapitelkopfzeile}
\textbf{\LaTeX-Vorlage für Projekt-, Seminar- und Bachelorarbeiten}

Bei dem vorliegenden Dokument handelt es sich um eine Vorlage, die
für Projekt-, Seminar- und Bachelorarbeiten im Studiengang
Wirtschaftsinformatik der DHBW Stuttgart verwendet werden kann.

Sie setzt die technischen Vorgaben der Zitierrichtlinien\footnote{Sie finden diese unter \enquote{Prüfungsleistungen} im Studierendenportal (\url{https://studium.dhbw-stuttgart.de/winf/pruefungsleistungen/}).} des Studiengangs
(Stand: 01/2020) um.

\emph{Hinweise:} Bitte lesen Sie sich die Zitierrichtlinien unbedingt genau durch. Dieses Dokument ersetzt keine Anleitung oder Einführung in \LaTeX,
für die Nutzung sind daher gewisse Vorkenntnisse unerlässlich. Ein Einstieg in 
\LaTeX\ ist aber weniger schwierig, als es vielleicht auf den ersten Blick scheint
und lohnt sich für das Verfassen wissenschaftlicher Arbeiten in jedem Fall.\footnote{%
so auch \url{http://www.spiegel.de/netzwelt/tech/textsatz-keine-angst-vor-latex-a-549509.html}} 
Als Hilfestellung beim Schreiben eines Dokuments habe ich einen zweiseitigen kompakten \LaTeX-Spickzettel erstellt, der über Moodle verfügbar ist.

Ihre Rückmeldungen und Anregungen zu dieser Vorlage nehme ich gerne per E-Mail an die Adresse
\url{tobias.straub@dhbw-stuttgart.de} entgegen.

--- Prof. Dr. Tobias Straub

\vspace{5em}

\begin{center}\small
\begin{tabular}{ccl}
\multicolumn{3}{c}{\textbf{Versionshistorie}}\\
\hline
1.0	& 05.02.2015 & erste Fassung \\
\hline
1.1 & 16.02.2015 & siehe~\ref{anhang:ReleaseNotes11} \\
\hline
1.2 & 20.04.2015 & siehe~\ref{anhang:ReleaseNotes12} \\
\hline
1.3 & 20.02.2016 & siehe~\ref{anhang:ReleaseNotes13} \\
\hline
1.4 & 24.07.2017 & siehe~\ref{anhang:ReleaseNotes14} \\
\hline
1.5 & 07.01.2018 & siehe~\ref{anhang:ReleaseNotes15} \\
\hline
1.6 & 07.04.2018 & siehe~\ref{anhang:ReleaseNotes16} \\
\hline
1.7 & 12.02.2019 & siehe~\ref{anhang:ReleaseNotes17} \\
\hline
1.8 & 10.02.2020 & siehe~\ref{anhang:ReleaseNotes18} \\
\end{tabular}
\end{center}

\end{abstract}


\cleardoublepage

%%% Inhalts-, Abbildungs-, Tabellenverzeichnisse %%%
% sollen einzeilig gesetzt werden, um Platz zu sparen 
\begin{spacing}{1}
\tableofcontents
\clearpage
\chapter*{Abkürzungsverzeichnis}
\addcontentsline{toc}{chapter}{Abkürzungsverzeichnis}

Ein Abkürzungsverzeichnis ist optional. Das Paket \verb|acronym| kann weit mehr, als hier gezeigt.\footnote{siehe \url{http://ctan.org/pkg/acronym}}
Beachten Sie allerdings, dass Sie die Einträge selbst in sortierter Reihenfolge angeben müssen.

\begin{acronym}[DHBW] 
% Argument definiert die Breite der ersten Spalte anhand des längsten vorkommenden Eintrags
\acro{CRM}{Customer Relationship Management}
\acro{DHBW}{Duale Hochschule Baden-Württemberg}
\acro{IEEE}{Institute of Electrical and Electronics Engineers}
\acro{ITIL}{IT Infrastructure Library}
\acro{RoI}{Return-On-Invest}
\acro{UCS}{Universal Character Set}
\acro{UTF-8}{8-Bit UCS Transformation Format} 
\end{acronym}

\vspace{2em}
{\footnotesize
\textbf{Ergänzende Bemerkung:}
Eine im Text verwendete Abkürzung sollte bei ihrer ersten Verwendung erklärt werden. Falls Sie sich nicht selbst darum kümmern möchten, kann das das Paket \verb|acronym| übernehmen und auch automatisch Links zum Abkürzungsverzeichnis hinzufügen. Dazu ist an allen Stellen, an denen die Abkürzung vorkommt, \verb|\ac{ITIL}| zu schreiben. 

Das Ergebnis sieht wie folgt aus: 
\begin{itemize}
\item erstmalige Verwendung von \verb|\ac{ITIL}| ergibt: \ac{ITIL},
\item weitere Verwendung von \verb|\ac{ITIL}| ergibt: \ac{ITIL}
\end{itemize}
Wo benötigt, kann man mit dem Befehl \verb|\acl{ITIL}| wieder die Langfassung ausgeben lassen: \acl{ITIL}.

Falls man die Abkürzungen durchgängig so handhabt, kann man durch Paket-Optionen (in \verb|_dhbw_praeambel.tex|)
erreichen, dass im Abkürzungsverzeichnis nur die tatsächlich verwendeten Quellen aufgeführt werden (Option: \verb|printonlyused|) und zu jedem Eintrag die Seite der ersten Verwendung angegeben wird (Option: \verb|withpage|).
}

\clearpage
\thispagestyle{kapitelkopfzeile}
\listoffigures
\phantomsection
\addcontentsline{toc}{chapter}{Abbildungsverzeichnis} % Abb.verz. ins Inh.verz. aufnehmen

\clearpage
\listoftables
\phantomsection
\addcontentsline{toc}{chapter}{Tabellenverzeichnis}   % Tab.verz. ins Inh.verz. aufnehmen
\end{spacing}

%%% Umstellung der Seiten-Nummerierung auf 1, 2, 3 ... %%%
\cleardoublepage
\pagenumbering{arabic}

%%% Ihr eigentlicher Inhalt %%%
% Empfehlung: strukturieren Sie Ihren Text in einzelnen Dateien 
% und binden Sie diese hier mit \input{includes/dateiname.tex} ein

\chapter{Einleitung}

Bald kann nun der Text Ihrer Projekt- oder Bachelorarbeit beginnen. Dank \LaTeX\ wird Ihre Arbeit garantiert professionell aussehen. Für den Inhalt sind Sie aber weiterhin selbst verantwortlich~\mbox{;-)}

Natürlich ist es schwer, sich vorzustellen, wie das Dokument aussieht, wenn die Vorlage doch gar keinen Text enthält. Aus diesem Grund wird mit Hilfe des Pakets \enquote{blindtext} so genannter Blindtext erzeugt. Mit dem Befehl \verb|\blinddocument| wird nachfolgend ein ganzes Kapitel sinnfreier Blindtext eingefügt.\footnote{Beachten Sie, dass Sie in Ihrer Arbeit eine Strukturierung wie in Abschnitt 2.1 vermeiden sollten: Dort gibt es einen Abschnitt 2.1.1, aber keinen Abschnitt 2.1.2.} 

In den Abschnitten~\ref{section:werkzeuge} bis \ref{section:zeichencodierung} sind zuvor die wichtigsten Werkzeuge, die Dateistruktur der Vorlage sowie einige Einstellungen beschrieben. Abschnitt~\ref{section:fehlerbehebung} gibt Hilfestellungen für bestimmte Fehler. In Kapitel~\ref{chapter:zitate} finden sich Beispiele, wie Sie Quellen korrekt zitieren können. In Kapitel~\ref{chapter:abbildungenTabellen} werden Abbildungen, Tabellen, ein Code-Listing und auch mathematische Formeln in den Text eingebunden. Ab Seite~\pageref{chapter:quellen} finden Sie das Literaturverzeichnis.

\emph{Hinweis:} Die farbigen, anklickbaren Links, die in der PDF-Ansicht enthalten sind, werden beim Ausdruck nicht wiedergegeben.\footnote{Das Feature lässt sich abschalten, indem man die Option \texttt{hidelinks} bei \texttt{documentclass} zu Beginn des Hauptdokuments hinzufügt. Für weitere Konfigurationsmöglichkeiten siehe \url{https://ctan.org/pkg/hyperref}.}

\section{Werkzeuge}\label{section:werkzeuge}

Sämtliche benötigten Werkzeuge sind Open Source und damit kostenlos nutzbar.

Für einen einfachen Start in \LaTeX\ sowie Tests können Sie einen Online-Editor wie \url{overleaf.com} verwenden. Allerdings sollten Sie eine Projekt- oder Bachelorarbeit nicht mit diesem System erstellen, da es sich um eine Cloud-Lösung handelt, bei der letztlich nicht gewährleistet ist, dass Firmen-Interna und schützenswerte Daten nicht in die Hände Dritter gelangen können. Für eine lokale Installation können Sie z.B.\ TeXLive (alle gängigen Plattformen, \url{http://tug.org/texlive/}), MikTeX (Windows, \url{http://www.miktex.org/}) oder TexShop (Mac OS, \url{http://pages.uoregon.edu/koch/texshop/}) verwenden.

In Moodle findet sich der Link zu einer portablen Version von MikTeX, welche ohne Installation auskommt. Dieses System wird für Schulungen an der DHBW genutzt, es enthält aber prinzipiell alle Werkzeuge, um eine Bachelorarbeit zu schreiben.\footnote{Manche werden das relativ schlichte TeXworks durch einen anderen Editor ersetzen wollen.} Mit diesem System wird auch jede neue Version der Vorlage getestet.

Erstellen Sie Ihr Dokument von Beginn an in \LaTeX. Es ist etwa wenig sinnvoll, zuerst in Word zu schreiben und das Ergebnis am Ende nach \LaTeX\ zu konvertieren.\footnote{Falls es je nötig sein sollte, gibt es für die umgekehrte Richtung diverse Konverter, z.B.\ latex2rtf, \url{http://sourceforge.net/projects/latex2rtf/}.} 

   
Die Vorlage verwendet biblatex,\footnote{siehe etwa \url{http://www.ub.uni-konstanz.de/serviceangebote/literaturverwaltung/bibtex/bibtex-und-biblatex-benutzen.html} für eine Gegenüberstellung von BibTeX und biblatex.}
um den vorgegebenen Stil der Zitierrichtlinien umzusetzen. Der Editor muss daher so konfiguriert werden, dass er beim Übersetzen biber statt bibtex verwendet. Zur Verwaltung der Quellen nutzt man etwa JabRef (Java, \url{http://jabref.sourceforge.net/}).

\section{Benötigte Dateien}

Die folgenden Dateien, deren Namen einheitlich mit \verb|_dhbw_| beginnt, werden eingebunden. Änderungen an diesen Dateien sind nicht erforderlich.

\verb|_dhbw_authoryear.bbx       (Anpassung Einträge im Literaturverzeichnis)| \\
muss sich im selben Verzeichnis wie die zu kompilierende Hauptdatei (\verb|latex-vorlage.tex|) befinden, die folgenden Dateien im Unterverzeichnis \verb|template|:
\begin{verbatim}
_dhbw_biblatex-config.tex  (weitere Einstellung für Biblatex)
_dhbw_erklaerung.tex       (ehrenwörtliche Erklärung)
_dhbw_kopfzeilen.tex       (Kapitelname in Kopfzeilen) 
_dhbw_praeambel.tex        (Einbindung der benötigten Pakete)
\end{verbatim}

Weiterhin sind nur für die Übersetzung der Beispieldatei erforderlich:
\begin{verbatim}
dhbw.png                   (Beispiel für eine Grafik, die eingebunden wird)
HelloWorld.java            (eingebundenes Java-Listing)
literatur-datenbank.bib    (Literatur-Datenbank mit Beispiel-Einträgen)
\end{verbatim}
Die Grafik ist im Verzeichnis \verb|graphics| abgelegt, die anderen Dateien in \verb|includes|. 

Ferner gliedert sich dieser Text in folgende Dateien (ebenfalls im Verzeichnis \verb|includes|), die per \verb|\input|-Befehl eingebunden werden:
\begin{verbatim}
abbildungen_und_tabellen.tex
abkuerzungen.tex
abstract.tex
anhang.tex
deckblatt.tex
einleitung.tex
text_mit_zitaten.tex
\end{verbatim}

Um das Editieren (und Debuggen) zu erleichtern, ist es ratsam, ein längeres Dokument in einzelne Dateien zu strukturieren (z.B.\ kapitelweise) und ggf.\ die Ordnerstruktur nach eigenen Bedürfnissen anzupassen. Viele Editoren unterstützen es, in einem Projekt die Hauptdatei festzulegen, so dass der Compiler von einer beliebigen Unterdatei aus aufgerufen werden kann. (In TeXworks geschieht dies durch die Angabe \verb|% !TeX root = ../latex-vorlage.tex| in der ersten Zeile einer Unterdatei, wobei die Hauptdatei relativ zu dieser angegeben werden muss.)

\section{Vorzunehmende Einstellungen}

\emph{Hinweis:} Mit der ab 1/2016 geltenden Fassung der Zitierrichtlinien ist die zuvor bestehende Möglichkeit, auch doppelseitig auszudrucken, gestrichen worden. Gleiches gilt für die nun nicht mehr bestehende Option, \enquote{ebenda} zu nutzen.

\subsection{Deckblatt}

Die folgenden Anpassungen sind in der Datei \verb|deckblatt.tex| vorzunehmen. Die entsprechenden Stellen sind im Source Code wie folgt gekennzeichnet:
\lstset{language=TeX} 
\begin{lstlisting}
% HIER EDITIEREN: 
\end{lstlisting}

\subsubsection{Titel und Erklärung}
Natürlich ist das Deckblatt anzupassen, schließlich soll dort Ihr Name erscheinen. Denken Sie bitte aber auch unbedingt daran, am Anfang der Datei festzulegen, um welchen Typ (Projekt-/ Bachelor-/Seminararbeit) es sich handelt und wie Ihr Thema (Titel der Arbeit) lautet.\footnote{%
Falls Sie einen langen Titel mit Untertitel haben, können Sie für die Aufnahme in die Erklärung den Untertitel auch weglassen.} Beide Angaben werden automatisch in die ehrenwörtliche Erklärung eingefügt, so dass Sie in der Datei \verb|_dhbw_erklaerung.tex| keine Änderungen vornehmen müssen.
In der Datei \verb|deckblatt.tex| wird die Schriftgröße für den Titel auf \verb|\LARGE| eingestellt. Falls Sie je einen sehr langen Titel haben, wählen Sie eine kleinere Schriftgröße mit dem Befehl \verb|\Large| (Zeile 27 im Source Code).

\subsubsection{Meta-Daten im PDF}
Zu Beginn von \verb|deckblatt.tex| können Sie Ihren Namen und den Titel der Arbeit für die Meta-Daten der PDF-Datei angeben. 


\subsubsection{Vertraulichkeitsvermerk}
Ein so genannter \enquote{Sperr- oder Vertraulichkeitsvermerk} sollte eher die Ausnahme sein. Meiner Erfahrung nach enthalten die allerwenigsten Arbeiten brisante Firmengeheimnisse. Arbeiten werden -- unabhängig von einem Sperrvermerk -- von der DHBW sowieso nicht an Dritte weitergegeben, auch die Prüfer müssen die Inhalte vertraulich behandeln. Durch einen Sperrvermerk schränken Sie sich aber möglicherweise selbst ein, da Sie dann Ihre Arbeit im Grunde auch niemandem ohne Genehmigung Ihrer Firma zeigen dürfen.
Falls ein Sperrvermerk erforderlich ist, können Sie die auskommentierten Zeilen am Ende der Datei \verb|deckblatt.tex| nutzen.

\subsection{Hauptdokument}

Sie bearbeiten folgenden Abschnitt, um Ihre einzelnen Kapitel einzubinden:
\lstset{language=}
\lstinputlisting[firstline=48,lastline=59,firstnumber=48]{latex-vorlage.tex}

Je nachdem, ob Sie Abstract und Abkürzungsverzeichnis verwenden, ändern Sie die folgenden Zeilen.
In aller Regel werden Sie auch Tabellen und Abbildungen nutzen, ansonsten entfernen Sie die entsprechenden Verzeichnisse. 
\lstinputlisting[firstline=21,lastline=42,firstnumber=21]{latex-vorlage.tex}

Änderungen außerhalb dieser Bereiche sind nicht erforderlich.

\emph{Hinweis:} Mit den Änderungen der Zitierrichtlinien 01/2020 ist die Anforderung entfallen, ein Gesprächsverzeichnis im Quellenverzeichnis zu führen. Stattdessen ist dies in den Textteil bzw.\ einen Anhang zu integrieren.


\section{Zeichencodierung}\label{section:zeichencodierung}

Achten Sie darauf, dass Sie für alle \LaTeX\ und BibTeX-Dateien eine einheitliche Zeichencodierung verwenden, damit Umlaute und Sonderzeichen korrekt wiedergegeben werden. Sie müssen dazu ggf.\ die Einstellungen Ihres Editors anpassen.

Damit Sonderzeichen korrekt dargestellt werden, sollte als Codierung \emph{UTF-8} (Unicode) eingestellt sein. UTF-8 stimmt in den ersten 128 Unicode-Zeichen mit dem ASCII-Zeichensatz überein, kann aber auch Sonderzeichen oder Zeichen beliebiger Sprachen darstellen.

Diese Datei ist \emph{UTF-8}-codiert, weshalb sich in der Präambel folgender Befehl findet:
\lstset{language=TeX} 
\begin{lstlisting}
\usepackage[utf8]{inputenc}   % Zeichencodierung UTF-8 für Eingabe-Dateien
\end{lstlisting}

Mit der richtigen Codierung können die Sonderzeichen ä, ö, ü, Ä, Ö, Ü, ß, € wie gewohnt direkt im Source Code eines \LaTeX-Dokuments geschrieben werden.

\section{Fehlerbehebung}\label{section:fehlerbehebung}

\subsection{Kopfzeile}
Es ist \emph{kein Fehler}, sondern beabsichtigt, dass auf der ersten Seite eines Kapitels (und auch bei der ersten Seite der Verzeichnisse) die Kopfzeile fehlt.

\subsection{Kontrolle der Seitenränder}
Wenn Sie in Zeile 3 der Hauptdatei das Packet \enquote{showframe} einbinden, können Sie sich die Ränder Ihres Dokuments zur Kontrolle anzeigen lassen. Es empfiehlt sich, die Ränder zu kontrollieren und darauf zu achten, dass beim Druck \emph{keine Skalierung} im PDF-Viewer gewählt wird.  

Überlange Zeilen erkennen Sie außerdem an den Warnmeldungen \verb|Overfull \hbox| in der \verb|.log|-Datei.

\subsection{biber}
Zumindest unter Mac OS kommt es manchmal vor, dass sich biber aufhängt und mit einer Fehlermeldung folgender Form stoppt:

{\small
\begin{verbatim}
read_file '/var/folders/ay/ay9RQK7FEcKOxehY75+N4k+++TI/-Tmp-/
par-746f62696173737472617562/cache-a3cdad92316c60c9c5179d80d6bb51a7a024393c/
inc/lib/Biber/LaTeX/recode_data.xml' - sysopen: No such file or directory at 
/var/folders/ay/ay9RQK7FEcKOxehY75+N4k+++TI/-Tmp-/par-746f62696173737472617562/
cache-a3cdad92316c60c9c5179d80d6bb51a7a024393c/inc/lib/Biber/LaTeX/Recode.pm 
line 112.
INFO - This is Biber 1.9
INFO - Logfile is 'latex-vorlage.blg'
\end{verbatim}
}

Abhilfe schafft Löschen des temporären Verzeichnisses, hier: \verb|par-746f62696173737472617562|, inklusive der Unterverzeichnisse.

Ein korrekter biber-Lauf sieht in etwa so aus:
{\small
\begin{verbatim}
INFO - Logfile is 'latex-vorlage.blg'
INFO - Reading 'latex-vorlage.bcf'
INFO - Found 25 citekeys in bib section 0
INFO - Processing section 0
INFO - Looking for bibtex format file 'includes/literatur-datenbank.bib' for section 0
INFO - Decoding LaTeX character macros into UTF-8
INFO - Found BibTeX data source 'includes/literatur-datenbank.bib'
INFO - Overriding locale 'de-DE' defaults 'variable = shifted' with 'variable = 
non-ignorable'
INFO - Overriding locale 'de-DE' defaults 'normalization = NFD' with 
'normalization = prenormalized'
INFO - Sorting list 'nyt' of type 'entry' with scheme 'nyt' and locale 'de-DE'
INFO - No sort tailoring available for locale 'de-DE'
INFO - Writing 'latex-vorlage.bbl' with encoding 'UTF-8'
INFO - Output to latex-vorlage.bbl 
\end{verbatim}
}

\subsection{Zeilenumbruch bei langen URLs}
Sehr lange und komplizierte URLs\footcite{langeURL} können in manchen Fällen ein Problem beim Zeilenumbruch in Blocksatz darstellen, v.a.\ im Verzeichnis der Internetquellen. Folgende URL ragt beispielsweise in den rechten Rand hinein: 

\url{http://www.google.de/search?hl=de&source=hp&q=biblatex+umbruch+url&gbv=2&oq=biblatex+umbruch+&gs_l=heirloom-hp.3.0.0i13i30l2j0i22i10i30.1757.7464.0.8525.22.19.0.3.3.0.204.1959.13j5j1.19.0.msedr...0...1ac.1.34.heirloom-hp..0.22.1990.yWySyKFfLPY}

Eine nahe liegende Lösungsmöglichkeit ist, den betroffenen Satz etwas umzuformulieren.
Jetzt geht es allerdings auch nicht besser.\footnote{%
\LaTeX\ sieht von einer Trennung der URL bei einem Bindestrich ab, um Missverständnisse zu vermeiden.
}
\url{http://www.google.de/search?hl=de&source=hp&q=biblatex+umbruch+url&gbv=2&oq=biblatex+umbruch+&gs_l=heirloom-hp.3.0.0i13i30l2j0i22i10i30.1757.7464.0.8525.22.19.0.3.3.0.204.1959.13j5j1.19.0.msedr...0...1ac.1.34.heirloom-hp..0.22.1990.yWySyKFfLPY}

Man kann sich aber damit behelfen, dass man die URL manuell in mehrere hintereinander gestellte \verb|\href|-Befehle auftrennt (der anklickbare Link bleibt so weiterhin korrekt). Einzig die Zeichen \verb|&| und \verb|_| sind zu escapen (als \verb|\&| bzw.\ \verb|\_|). Schauen Sie bitte in den Source Code.

\newcommand{\mylongurl}[1]{%
% erstes Argument von href ist das Linkziel, das beim Klicken angesprungen wird
\href{http://www.google.de/search?hl=de&source=hp&q=biblatex+umbruch+url&gbv=2&oq=biblatex+umbruch+&gs_l=heirloom-hp.3.0.0i13i30l2j0i22i10i30.1757.7464.0.8525.22.19.0.3.3.0.204.1959.13j5j1.19.0.msedr...0...1ac.1.34.heirloom-hp..0.22.1990.yWySyKFfLPY}{\texttt{#1}}\\}

%beachten Sie, dass & und _ escaped werden müssen
\mylongurl{http://www.google.de/search?hl=de\&source=hp\&q=biblatex+umbruch+url\&gbv=2\&oq=}
\mylongurl{biblatex+umbruch+\&gs\_l=heirloom-hp.3.0.0i13i30l2j0i22i10i30.1757.7464.0.8525.22.}
\mylongurl{19.0.3.3.0.204.1959.13j5j1.19.0.msedr...0...1ac.1.34.heirloom-hp..0.22.1990.}
\mylongurl{yWySyKFfLPY}

\subsection{Verwendung von MikTeX Portable}
Falls Sie MikTeX Portable verwenden, nutzen Sie bitte die Funktion \enquote{Check for updates}, welche über das Tray Icon erreichbar ist. In einem Fall gab es vor dem Update Schwierigkeiten mit der automatischen Silbentrennung des Babel-Pakets, die sich im Compiler-Log wie folgt äußerten
{\small\begin{verbatim}
Package babel Warning: No hyphenation patterns were preloaded for
(babel)                the language `German (new orthography)' into the format.

(babel)                Please, configure your TeX system to add them and
(babel)                rebuild the format. Now I will use the patterns
(babel)                preloaded for \language=0 instead on input line 43.
\end{verbatim}
}
Bei korrekter Funktionsweise sollte sich im Log stattdessen nur ein Hinweis finden wie:
{\small\begin{verbatim}
Babel <3.18> and hyphenation patterns for 75 language(s) loaded.
\end{verbatim}
}
\blinddocument
\chapter{Zitieren}\label{chapter:zitate}

Der Zitierstil ist so angepasst, dass er den Zitierrichtlinien des Studiengangs Wirtschaftsinformatik der DHBW Stuttgart entspricht. 

\section{Zitate in den Text einfügen}
In \LaTeX\ wird mit den Befehlen \verb|\footcite| 
oder 
\verb|\cite|
eine Referenz im Text eingefügt. Meist wird \verb|\cite| nur \emph{innerhalb} einer Fußnote benutzt. 
Damit ein vorangestelltes \enquote{Vgl.} in der Fußnote erscheint, können Sie wie folgt zitieren:
\begin{verbatim}
\footcite[Vgl.][S. 3]{Autor}
\footcite[Vgl.][]{Autor}
\end{verbatim}

Das erste optionale Argument von \verb|\footcite| wird dem Zitat vorangestellt, das zweite ist die Seitenzahl. Den selben Effekt hätte
\begin{verbatim}
\footnote{Vgl. \cite[S. 3]{Autor}}
\footnote{Vgl. \cite{Autor}}
\end{verbatim}

Hinweis: Falls \enquote{Vgl.}, aber keine Seitenzahl angeben werden soll, muss das zweite Argument vorhanden (jedoch leer) sein, ansonsten wird \enquote{Vgl.} als Seitenzahl interpretiert. Falsch ist also: 
\begin{verbatim}
\footcite[Vgl.]{Autor}  % so nicht!
\end{verbatim}



\subsection{Beispiele}
Nachfolgend ein paar Beispiele, um die korrekte Darstellung zu überprüfen:

\begin{itemize}
\item \cite{Schlosser} ist ein Buch über \LaTeX.
\item Zur Vorlesung \emph{Logik und Algebra} gibt es das gleichnamige Lehrbuch.\footcite{Staab}
\item nochmal dasselbe Buch\footcite{Staab}
\item ein weiteres Buch desselben Autors\footcite{BuschlingerStaab}

\item Der Konferenzbeitrag \cite{Ancuti} beschäftigt sich mit Bildverarbeitung.
\item Cloud Computing wird in einer Diplomarbeit erklärt.\footcite[S.~14]{Boettger:Diplomarbeit}
\item Preiß\footcite{Preiss} gibt eine Einführung in Datenbanken.
\item Eine Erläuterung, was \enquote{Intangibles} sind, findet sich bei Stoi\footcite[S.~82]{Stoi}.
\item weitere Ausführung in derselben Quelle\footcite[Vgl.][S.~84]{Stoi}
\item Laut Wikipedia\footcite[Abschnitt~5]{wiki:Wirtschaftsinformatik} ist Wirtschaftsinformatik ein interessantes Studienfach.
\item ITIL-Prozesse kann man tatsächlich auch mit \LaTeX\ dokumentieren.\footcite{Carvalho:PJ:2012-1}
\item Open-Source und Cloud-Computing in einem Buchbeitrag\footcite{Wind}
\item Buch mit zwei Autoren\footcite{MitZweiAutoren}
\item Buch mit drei Autoren\footcite{MitDreiAutoren}
\item Buch ohne Autor\footcite{OhneAutoren}
\item Buch ohne Autor und ohne Jahr\footcite{OhneAutorenOhneJahr}
\item und noch ein anderes Buch ohne Autor und ohne Jahr\footcite{OhneAutorenOhneJahr2}
\item Buch ohne Autor, aber dafür mit Herausgeber\footcite{keinAutorAberHerausgeber}

\item manche Bachelorarbeit baut auf einer vorhergehenden Projektarbeit\footcite{mayer:PA1} auf
\item das Handbuch zu BibLaTeX\footcite{biblatex:manual} und eines zu Windows 8\footcite{Win8}
\item zwei Beiträge zu Büchern\footcite{Trautwein:Nokia}${}^{,}$\footcite{Mann} und zu einem Konferenzband\footcite{Trautwein:Erfolgsfaktoren}
\item eine Online-Quelle\footcite{SAP:HANA}
\item eine plagiierte Dissertation,\footcite{GuttenPlag} nicht zur Nachahmung empfohlen
\item zum Testen, ob Umlaute und Sonderzeichen korrekt wiedergegeben werden\footcite{Umlauttest}
\end{itemize}

\subsection{Spezialfälle}
\begin{itemize}
\item \emph{Zwei Quellen am Satzende} werden durch Komma getrennt.\footcite{Staab}${}^{,}$\footcite{mayerLukas:PA1} Hier muss \verb|${}^{,}$| eingeschoben werden.
\item \emph{Eindeutigkeit:} Normalerweise wird kein Vorname des Autors angegeben. Falls es allerdings zur Eindeutigkeit\footnote{\cite{trautwein2011unternehmensplanspiele} vs. \cite{hitzler2011optimierung}} (bei gleicher Jahreszahl) erforderlich ist, wird der Vorname abgekürzt bzw.\ nötigenfalls sogar ganz ausgeschrieben mit angegeben.\footnote{Vgl. \cite{mayer:PA1} und \cite{mayerLukas:PA1}}
 
Welch ein Glück, dass Sie sich darum dank \LaTeX\ gar nicht kümmern müssen (arme Word\texttrademark-User ;-).

\item Die Verwendung von \emph{Sekundärliteratur}\footcitePrimaerSekundaer{Primaerquelle}{}{Sekundaerquelle}{11} wird weiter in Abschnitt~\ref{sec:sekundaerliteratur} erläutert.
\end{itemize}


\section{Eintragstypen für die Literatur-Datenbank}

Die verwendete Literatur pflegen Sie in einer Literatur-Datenbank im Bibtex-Format. Dabei handelt es sich um eine Textdatei, wobei für jede Quelle mittels Name-Value-Pairs die relevanten Attribute (Autor, Titel etc.) hinterlegt sind. Die Datei wird üblicherweise nicht im Texteditor, sondern in einem spezialisierten Programm wie JabRef bearbeitet.

Sofern in der Literatur-Datenbank der Typ eines Eintrags (Entry Type) korrekt festgelegt ist, wird er im Literaturverzeichnis automatisch richtig dargestellt. Mit folgenden Typen sollten Sie i.d.R.\ auskommen:
\begin{description}
\item[article] Artikel in einer Fachzeitschrift, auch E-Journal (Zeitschrift in elektronischer Form)\footnote{Bei E-Journals/E-Books werden beim Zitieren anstelle der (u.U. nicht eindeutigen, da von der Schriftgröße abhängigen) Seitenzahl Abschnitt und Absatz näher bezeichnet, also: \cite[Abschnitt 1.2.3, Absatz 4]{Staab}.}
\item[book] Buch, auch E-Book 
\item[inbook] Kapitel in einem Buch, zu dem mehrere Autoren beigetragen haben 
\item[inproceedings] Beitrag zu einer Fachtagung/Konferenz 
\item[manual] Handbuch
\item[misc] anderweitig nicht zuordenbarer Typ
\item[phdthesis] Dissertation
\item[thesis] Bachelor-/Master-/Diplomarbeit (Art wird im Attribut \enquote{type} festgelegt) 
\item[online] Internet- oder Intranet-Quelle\footnote{\label{fn:onlineEntryType}Man beachte, dass der Eintragstyp \enquote{online} in JabRef nur im \enquote{biblatex-Modus} (Menü: Datei -- Neue biblatex Bibliothek) auswählbar ist.}
\item[report] technischer Bericht, Forschungsbericht oder White Paper; diesen Typ können Sie auch verwenden, um eine Projektarbeit zu zitieren (Art wird im Attribut \enquote{type} festgelegt) 
\end{description} 

Eine Übersicht über die notwendigen Attribute jedes Eintragstyps gibt die folgende Tabelle, wobei ein Schrägstrich als \enquote{oder} zu verstehen ist.\footnote{Auszugsweise entnommen aus \cite{biblatex:manual}.} Zudem sind die wichtigsten optionalen Attribute aufgeführt.

\begin{tabular}{l|p{6cm}|p{6cm}}
\textbf{Eintragstyp} & \textbf{notwendige Attribute} & \textbf{optionale Attribute (Auswahl)} \\
\hline
article & author, title, journal, year/date  & volume, number, pages, month, note \\
\hline
book & author, title, year/date  & publisher, edition, editor, \mbox{volume/number}, series, isbn, url \\
\hline
inbook & author, title, booktitle, year/date  & bookauthor, editor, volume/number, series, isbn, url \\
\hline
inproceedings & author, title, booktitle, year/date  & organization/publisher, editor, volume/number, series, isbn, url  \\
\hline
manual & author/editor, title, year/date  & organization/publisher, address, edition, month, note, url, urldate \\
\hline
misc & author/editor, title, year/date & howpublished, organization, month, note \\
\hline
phdthesis & author, title, institution, year/date  & address, month, note \\
\hline
thesis & author, title, institution, type, year/date &  address, month, note \\
\hline
online\textsuperscript{\ref{fn:onlineEntryType}} & author/editor, title, year\footnotemark/date, url & urldate \\
\hline
report & author, title, institution, type, year/date & number, version, url, urldate \\
\end{tabular}
\footnotetext{%
Sofern kein Jahr bekannt ist, sollte das Attribut nicht leer gelassen werden (sonst wird die aktuelle Jahreszahl automatisch eingefügt), sondern der Eintrag \enquote{o.J.} gewählt werden.}



\section{Zitieren von Sekundärliteratur}\label{sec:sekundaerliteratur}
Gelegentlich lässt es sich nicht vermeiden, aus der Sekundärliteratur zu zitieren. Dies leistet der folgende Befehl.
\begin{verbatim}
\footcitePrimaerSekundaer{Primaerquelle}{Seite}{Sekundaerquelle}{Seite}
\end{verbatim}
Die erste Seitenangabe bezieht sich auf die Primär-, die zweite auf die Sekundärquelle.
Die Seitenangaben sind optional, sie können auch leer bleiben.\footcitePrimaerSekundaer{Primaerquelle}{23}{Sekundaerquelle}{}
Es ist aber zu beachten, dass der Befehl \verb|\footcitePrimaerSekundaer| vier Argumente hat.

Ins Literaturverzeichnis soll nur die Sekundärquelle aufgenommen werden. Dies wird dadurch erreicht, dass in der Literatur-Datenbank bei der Primärquelle im Attribut \enquote{keyword} der Wert \enquote{ausblenden} eintragen wird.

\chapter{Beispiele für Abbildungen und Tabellen}\label{chapter:abbildungenTabellen}

Hier finden Sie Beispiele für Abbildungen, Tabellen, Formelsatz und  Source Code.

\section{Abbildungen}
In diesem Abschnitt gibt die Abbildungen~\ref{abb:Logo2cmHoch} und~\ref{abb:Logo2cmBreit}, die beide das Logo der DHBW zeigen.

\begin{figure}[htb]
\centering
\includegraphics[height=2cm]{graphics/dhbw.png}
\caption[DHBW-Logo 2cm hoch]{DHBW-Logo 2cm hoch.\footnotemark}
\label{abb:Logo2cmHoch}
\end{figure}
\footnotetext{Mit Änderungen entnommen aus: \cite{OhneAutorenOhneJahr}}

\lstset{language=TeX, % hervorzuhebende Keywords definieren
  morekeywords={footnotetext,footnotemark,footcite,caption}
}

\emph{Spezialfall:} Sofern \emph{innerhalb} der Bezeichnung einer Abbildung eine Fußnote angegeben oder eine Quelle referenziert werden soll, geschieht dies nicht per \lstinline|\footnote| oder \lstinline
|\footcite|. Vielmehr sind die Befehle \lstinline|\footnotemark| und \lstinline|\footnotetext| zu verwenden und außerdem das optionale Argument für \lstinline|\caption| anzugeben (vgl.\ Source Code).

\begin{figure}[htb]
\centering
\includegraphics[width=2cm]{graphics/dhbw.png}
\caption[DHBW-Logo 2cm breit.]{DHBW-Logo 2cm breit. (Quelle: DHBW\footnotemark)}
\label{abb:Logo2cmBreit}
\end{figure}
\footnotetext{\url{www.dhbw.de}}



\section{Tabellen}

In diesem Abschnitt gibt es zwei Beispiel-Tabellen, nämlich auf Seite~\pageref{tab:BeispielTabelleKlein} und auf Seite~\pageref{tab:BeispielTabelleGroesser}.

\begin{table}[htb]
\centering
\begin{tabular}{lcr}
links & Mitte & rechts \\
\hline
Muster & Muster & Muster \\
\end{tabular}
\caption{Kleine Beispiel-Tabelle.}
\label{tab:BeispielTabelleKlein}
\end{table}

\begin{table}[htb]
\centering
\begin{tabular}{|l|l|c|l|r||l}
    \textbf{Spalte 1} & \textbf{Spalte 2} & \textbf{Spalte 3} & \textbf{Spalte 4} & \textbf{Spalte 5} & \textbf{Spalte 6} \\
    \hline
    a        & b          & c                & d        & e        & f        \\
    Test     & Test, Test & Test, Test, Test & ~        & ~        & ~        \\
    1        & 2          & 3                & 4        & 5        & 6        \\
\end{tabular}
\caption{Größere Beispiel-Tabelle.}
\label{tab:BeispielTabelleGroesser}
\end{table}

\section{Etwas Mathematik}

Eine abgesetzte Formel:
\[
  \int_a^b x^2 \: \mathrm{d} x = \frac{1}{3} (b^3 - a^3)
\]

Es ist $a^2+b^2 = c^2$ eine Formel im Text.

\section{Source Code}

Source Code-Blöcke können auf folgende Arten eingefügt werden:

\lstset{language=Java}

Direkt im \LaTeX-Source Code:
\begin{lstlisting}
if(1 > 0) {
  System.out.println("OK"); 
} else {
  System.out.println("merkwuerdig");
}
\end{lstlisting}

oder eingefügt aus einer externen Datei.
\lstinputlisting{includes/HelloWorld.java}

\chapter*{Anhang}
\addcontentsline{toc}{chapter}{Anhang}
\section*{Anhangverzeichnis}
\vspace{-8em}

% vor \listofanhang müssen Einrückungen angepasst werden
\abstaendeanhangverzeichnis

\listofanhang
\clearpage
\spezialkopfzeile{Anhang} % damit in der Kopfzeile das Wort "Anhang" angezeigt wird

\anhang{So funktioniert's}

\lstset{language=TeX, % hervorzuhebende Keywords definieren
  morekeywords={anhang, anhangteil}
}


Um den Anforderungen der Zitierrichtlinien nachzukommen, wird das Paket \verb|tocloft| verwendet. Jeder Anhang wird mit dem (neu definierten) Befehl \lstinline|\anhang{Bezeichnung}| begonnen, der insbesondere dafür sorgt, dass ein Eintrag im Anhangsverzeichnis erzeugt wird. Manchmal ist es wünschenswert, auch einen Anhang noch weiter zu unterteilen. Hierfür wurde der Befehl \lstinline|\anhangteil{Bezeichnung}| definiert.

In~\ref{anhang:abbildung} finden Sie eine bekannte Abbildung und etwas Source Code in~\ref{anhang:sourcecode}. 

\anhangteil{Wieder mal eine Abbildung}\label{anhang:abbildung}
\begin{figure}[htb]
\centering
\includegraphics[width=0.9\linewidth]{graphics/dhbw.png}
\caption{Mal wieder das DHBW-Logo.}
\end{figure}

\anhangteil{Etwas Source Code}\label{anhang:sourcecode}
\lstinputlisting{includes/HelloWorld.java}


\anhang{Release Notes}
\anhangteil{Änderungen in Version 1.1}\label{anhang:ReleaseNotes11}
In Version 1.1 sind einige Rückmeldungen, die nach der Einführungsvorlesung am 6.2.2015 oder nach Veröffentlichung der Vorlage in Moodle eingegangen sind, berücksichtigt worden. Korrekturen sind mit \enquote{(Fix)} gekennzeichnet. 

\begin{itemize}
\item \verb|latex-vorlage.tex|
\begin{itemize}
\item (Fix) Abkürzungsverzeichnis wird vor Abbildungsverzeichnis platziert
\item (Fix) Abbildungs- und Tabellenverzeichnis in Inhaltsverzeichnis aufgenommen
\item (Fix) Quellenverzeichnis wird nun ohne Kapitelnummer dargestellt

\item eingebundene Dateien in Unterverzeichnissen \verb|includes| bzw.\ \verb|graphics|
\item Beispiel-Anhang (Datei \verb|anhang.tex|) mit Erklärungen wurde eingebunden 
\end{itemize}

\item \verb|_dhbw_praeambel.tex|
\begin{itemize}
\item (Fix) das Paket hyperref wird nach biblatex eingebunden, um ein Problem mit der Verlinkung der Fußnoten im PDF zu beheben
\item (Fix) Fußnoten  gemäß der Richtlinien fortlaufend nummeriert und nicht pro Kapitel
\item Einstellungen hinzugefügt, um Anhangsverzeichnis zu ermöglichen
\item bessere Kompatibilität zwischen KOMA-Script (scrreprt) und anderen Paketen mittels scrhack
\end{itemize}

\item \verb|_dhbw_biblatex-config.tex|
\begin{itemize}
\item (Fix) keine Abschnittsnummern für einzelne Verzeichnisse im Quellenverzeichnis
\end{itemize}

\item \verb|abbildungen_und_tabellen.tex|
\begin{itemize}
\item Erklärung, wie eine Fußnote/ein Zitat bei einer Abbildung zu erstellen ist
\end{itemize}

\item \verb|abkuerzungen.tex|
\begin{itemize}
\item Abkürzungsverzeichnis wird im Inhaltsverzeichnis aufgeführt
\end{itemize}

\item \verb|abstract.tex|, \verb|anhang.tex|, \verb|einleitung.tex| 
\begin{itemize}
\item Erklärungen im Text ergänzt
\end{itemize}

\item \verb|deckblatt.tex|
\begin{itemize}
\item Meta-Daten (Autor, Titel) für die generierte PDF-Datei lassen sich nun festlegen
\end{itemize}

\end{itemize}


\anhangteil{Änderungen in Version 1.2}\label{anhang:ReleaseNotes12}
Über das Forum in Moodle sind einige Rückmeldungen eingegangen -- vielen Dank an alle, die dazu beigetragen haben. In der Version 1.2 wurden folgende Änderungen vorgenommen, wobei Korrekturen wieder mit \enquote{(Fix)} gekennzeichnet sind: 

\begin{itemize}
\item \verb|latex-vorlage.tex| (Hauptdokument)
\begin{itemize}
\item (Fix) Zeile 19: Seitenzahlen zu Beginn mit römischen \emph{Groß}buchstaben nummeriert
\end{itemize}

\item \verb|_dhbw_praeambel.tex|
\begin{itemize}
\item Zeile 39/40: Unterstützung für \enquote{ebenda} 
\item Zeile 46--68: zweite Gliederungsebene für Anhänge ermöglicht
\item (Fix) Zeile 70--73: Abbildungen und Tabellen: Zähler fortlaufend, kein Rücksetzen zu Kapitelbeginn (Paket \verb|chngcntr| anstelle von Paket \verb|remreset|)
\end{itemize}

\item \verb|_dhbw_biblatex-config.tex|
\begin{itemize}
\item (Fix) bei Quellen mit Herausgeber, aber ohne Autor wird der Name des Herausgebers im Verzeichnis fett gedruckt
\item Unterstützung für \enquote{ebenda} 
\end{itemize}

\item \verb|abkuerzungen.tex|
\begin{itemize}
\item Bemerkungen zur fortgeschrittenen Nutzung des \verb|acronym|-Pakets eingefügt 
\end{itemize}

\item \verb|einleitung.tex|
\begin{itemize}
\item Abschnitt 1.3 zu Einstellungen ergänzt
\item Abschnitt 1.5 zu Fehlerbehebungen eingefügt 
\end{itemize}

\item \verb|text-mit-zitaten.tex|
\begin{itemize}
\item Abschnitt 3.1 eingefügt, Erläuterungen zum Zitieren mit \enquote{vgl.} und \enquote{ebenda}. 
\item Abschnitt 3.2: Beispiele ergänzt
\item Hinweis zu Jahreszahlen bei Online-Quellen
\end{itemize}

\item \verb|anhang.tex|
\begin{itemize}
\item Erläuterungen zur zweiten Gliederungsebene
\end{itemize}

\item \verb|literatur-datenbank.bib|
\begin{itemize}
\item weitere Beispiele für Quellen
\end{itemize}

\end{itemize}

\anhangteil{Änderungen in Version 1.3}\label{anhang:ReleaseNotes13}
Durch die ab 1/2016 geltenden Änderungen der Zitierrichtlinien des Studiengangs waren einige kleinere Anpassungen der Vorlage erforderlich, die nachfolgend beschrieben sind. Bei dieser Gelegenheit ebenfalls erfolgte Korrekturen sind wieder mit \enquote{(Fix)} gekennzeichnet:

\begin{itemize}
\item \verb|latex-vorlage.tex| (Hauptdokument)
\begin{itemize}
\item Hinweis auf Option doppelseitiger Druck entfernt
\item Schriftgröße der Kapitelüberschriften verkleinert
\item (Fix) Kopf- und Fußzeilen werden nun korrekt angezeigt für erste Seite eines Kapitels und auch  Quellenverzeichnisse
\end{itemize}

\item \verb|_dhbw_praeambel.tex|
\begin{itemize}
\item Angabe des unteren Rands für Seitenzahl, da diese nun unten rechts steht
\item Unterstützung für \enquote{ebenda} entfernt
\item (Fix) Präfixe wie \enquote{von} im Namen eines Autors werden berücksichtigt
\item Anpassung der Abstände bei Kapitelüberschriften
\item Kopf- und Fußzeile für Verzeichnisse nun in \verb|_dhbw_kopfzeilen.tex| definiert 
\end{itemize}


\item \verb|deckblatt.tex|
\begin{itemize}
\item Schriftgröße des Titels vergrößert
\item Befehl \verb|\typMeinerArbeit| eingeführt, um Typ auszuwählen
\item Festlegung des Themas (für ehrenwörtliche Erklärung) mit Befehl \verb|\themaMeinerArbeit|
\item Darstellung der Angabe des Betreuers in der Ausbildungsstätte angepasst
\item Formulierung des Sperrvermerks angepasst  
\end{itemize}

\item \verb|_dhbw_erklaerung.tex|
\begin{itemize}
\item Formulierung angepasst an geänderte Prüfungsordnung
\item Typ und Thema der Arbeit werden automatisch eingefügt
\end{itemize}

\item \verb|_dhbw_kopfzeilen.tex|
\begin{itemize}
\item Seitennummern stehen jetzt unten rechts
\item (Fix) Kopf- und Fußzeile werden nun korrekt angezeigt in Verzeichnissen und dem Anhang
\end{itemize}

\item \verb|_dhbw_biblatex-config.tex|
\begin{itemize}
\item Anpassung des Zitierstils auf die ab 1/2016 geltenden Regelungen  
\item Vorkehrungen für Eindeutigkeit (Hinzufügen abgekürzter oder nötigenfalls ausgeschriebener Vorname) bei Übereinstimmung von Name und Jahreszahl 
\end{itemize}

\item \verb|einleitung.tex|
\begin{itemize}
\item Abschnitt 1.3 zu Einstellungen grundlegend überarbeitet
\item Abschnitt 1.5.2 zur Kontrolle der Seitenränder eingefügt 
\end{itemize}

\item \verb|text-mit-zitaten.tex|
\begin{itemize}
\item Abschnitt 3.1: Hinweise zu \enquote{ebenda} entfernt
\item Abschnitt 3.2: Beispiele zur Eindeutigkeit des Zitats ergänzt
\item Abschnitt 3.3: Hinweise für E-Journals/E-Books ergänzt 
\end{itemize}

\item \verb|anhang.tex|
\begin{itemize}
\item (Fix) Befehl \verb|\spezialkopfzeile| aufgenommen, damit in Kopfzeile das Wort \enquote{Anhang} angezeigt wird 
\item diese Release Notes wurden in eine eigene Datei verschoben
\end{itemize}

\item \verb|release_notes.tex|
\begin{itemize}
\item s.o.
\end{itemize}


\item \verb|literatur-datenbank.bib|
\begin{itemize}
\item weitere Beispiele für Quellen
\end{itemize}
\end{itemize}

\anhangteil{Änderungen in Version 1.4}\label{anhang:ReleaseNotes14}
Durch nicht abwärtskompatible Änderungen beim Versionswechsel von Biblatex 3.2 zu 3.3 sind einige Änderungen notwendig geworden.\footnote{Diese basieren auf Vorschlägen von Yannik Ehlert -- vielen Dank dafür!}
Die vorliegende Version 1.4 wurde erfolgreich mit MikTeX gestestet (portable Version 2.9.6361 vom 3.6.2017, unter Verwendung von Biblatex 3.7).

\begin{itemize}
\item \verb|_dhbw_biblatex-config.tex|
\begin{itemize}
\item Anpassung der \verb|\usebibmacro|-Befehle
\end{itemize}

\item \verb|_dhbw_authoryear.bbx|
\begin{itemize}
\item  Änderung von \verb|\printdateextralabel| zu \verb|\printlabeldateextra|
\end{itemize}
\end{itemize}

\anhangteil{Änderungen in Version 1.5}\label{anhang:ReleaseNotes15}
Für den Test dieser Version auf einem Windows-System wurde wieder die portable Version von MiKTeX (2.9.6521 vom 10.11.2017) verwendet.\footnote{\url{http://miktex.org/portable}} Da in diesem Paket leider die Versionen von Biblatex (3.10) und Biber (2.7) inkompatibel sind, ist es erforderlich, die Datei \verb|biber.exe| im Verzeichnis \verb|texmfs\install\miktex\bin\| durch die aktuelle Version 2.10 vom 20.12.2017\footnote{\url{https://sourceforge.net/projects/biblatex-biber/files/biblatex-biber/current/binaries/Windows/}} zu ersetzen. Im Editor TeXworks verwendet man dann zum Übersetzen des \LaTeX-Sourcecodes Typeset/pdfLaTeX bzw.\ Typeset/Biber.

Korrekturen sind wieder mit \enquote{(Fix)} gekennzeichnet.

\begin{itemize}
\item \verb|latex-vorlage.tex| (Hauptdokument)
\begin{itemize}
\item Nach der Änderung der Zitierrichtlinien gibt es nun kein separates Verzeichnis mehr für Internet- und Intranetquellen.
\item Option \verb|notkeyword=ausblenden| bei \verb|\printbibligraphy| sorgt dafür, dass Sekundärliteratur korrekt zitiert wird.
\end{itemize}

\item \verb|_dhbw_praembel.tex|
\begin{itemize}
\item (Fix) Die Bezeichnung geschachtelter Anhänge wurde auf das in den Zitierrichtlinien geforderte Format \enquote{Anhang 2/1} angepasst (Befehl \verb|\anhangteil|).
\end{itemize}

\item \verb|einleitung.tex|
\begin{itemize}
\item Hinweis zum Ausblenden der farbigen Links im PDF hinzugefügt
\end{itemize}

\item \verb|text-mit-zitaten.tex|
\begin{itemize}
\item Abschnitt 3.4 aktualisiert nach Wegfall des separaten Verzeichnisses für Internet- und Intranetquellen
\item Abschnitt zum Zitieren von Sekundärliteratur hinzugefügt
\end{itemize}

\end{itemize}


\anhangteil{Änderungen in Version 1.6}\label{anhang:ReleaseNotes16}
Diese Version wurde auf einem Windows-System erfolgreich mit der portablen Version von MiKTeX (2.9.6621 vom 18.02.2018) getestet.\footnote{Vielen Dank an Florian Eichin für seine wertvollen Anmerkungen.}

Korrekturen sind wieder mit \enquote{(Fix)} gekennzeichnet.

\newpage

\begin{itemize}
\item \verb|latex-vorlage.tex| (Hauptdokument)
\begin{itemize}
\item (Fix) An einer Stelle gab es in Version 1.5 (Internetquellen nicht mehr separat) noch ein Überbleibsel von Version 1.4 (Internetquellen separat), dies wurde korrigiert.
\item (Fix) Im Inhaltsverzeichnis war die Verlinkung des Abbildungs- und Tabellenverzeich\-nisses nicht ganz korrekt.
\item Mit den Befehlen \verb|\literaturverzeichnis| bzw.\ \verb|\literaturUndQuellenverzeichnis| kann bequem die Erstellung der Quellenverzeichnisse gesteuert werden, abhängig davon, ob es ein Gesprächsverzeichnis gibt oder nicht.
 
\end{itemize}

\item \verb|_dhbw_praembel.tex|
\begin{itemize}
\item Einrückungen für Abbildungs-, Tabellen- und Anhangverzeichnis angepasst
\item Abkürzungen \enquote{Abb.} und \enquote{Tab.} für Abbildungen bzw.\ Tabellen
\end{itemize}

\item \verb|_dhbw_biblatex-config.tex|
\begin{itemize}
\item Befehle \verb|\literaturverzeichnis| und \verb|\literaturUndGespraechsverzeichnis| definiert
\item Befehl \verb|\footcitePrimaerSekundaer| definiert
\end{itemize}

\item \verb|_dhbw_erklaerung.tex|
\begin{itemize}
\item Eintrag als \enquote{Erklärung} (statt \enquote{Ehrenwörtliche Erklärung}) ins Inhaltsverzeichnis
\end{itemize}

\item \verb|einleitung.tex|
\begin{itemize}
\item Bezeichnung \enquote{Erklärung} statt \enquote{Ehrenwörtliche Erklärung}
\item Erläuterung von \verb|\literaturverzeichnis| und \verb|\literaturUndGespraechsverzeichnis|
\item Hinweis auf Notwendigkeit von Updates bei MikTeX Portable
\end{itemize}

\item \verb|text_mit_zitaten.tex|
\begin{itemize}
\item Erläuterungen zu Befehl \verb|\footcitePrimaerSekundaer| ergänzt
\end{itemize}

\item \verb|anhang.tex|
\begin{itemize}
\item Befehl \verb|\abstaendeanhangverzeichnis| für Anpassung Einrückung ergänzt
\end{itemize}

\item \verb|literatur-datenbank.bib|
\begin{itemize}
\item Eintrag ergänzt
\end{itemize}

\end{itemize}

\anhangteil{Änderungen in Version 1.7}\label{anhang:ReleaseNotes17}
Diese Version wurde auf einem Windows-System erfolgreich mit der portablen Version von MiKTeX (2.9.6942 vom 04.01.2019) getestet.

Korrekturen sind wieder mit \enquote{(Fix)} gekennzeichnet.

\begin{itemize}
\item \verb|_dhbw-authoryear.bbx|
\begin{itemize}
\item Da \verb|labeldate| in Biblatex nicht mehr unterstützt wird, erfolgte eine Umbenennung in 
\verb|labeldateparts|.\footnote{vgl.\ \url{https://github.com/semprag/biblatex-sp-unified/issues/23}}
\end{itemize}

\item \verb|_dhbw_biblatex-config.tex|
\begin{itemize}
\item (Fix) Es wurde das Problem behoben, dass im Literaturverzeichnis bei bestimmten Eintragstypen der Titel in Anführungszeichen steht.\footnote{Danke an Florian Eichin für seinen Hinweis.}
\end{itemize}

\end{itemize}


\anhangteil{Änderungen in Version 1.8}\label{anhang:ReleaseNotes18}
Diese Version wurde auf einem Windows-System erfolgreich mit der portablen Version von MiKTeX (2.9.6942 vom 04.01.2019) getestet.

Die Aktualisierungen in der Vorlage spiegeln zum Einen die Änderungen in den Zitierrichtlinien wieder. Zum Anderen wurden einige studentische Vorschläge aufgegriffen, um die Nutzung der Vorlage zu erleichtern.\footnote{Danke an Bjarne Koll, Tobias Schwarz und Lars Ungerathen für ihre Anregungen.} 

\begin{itemize}

\item \verb|latex_vorlage.tex| (Hauptdokument)
\begin{itemize}
\item Es wird nun davon ausgegangen, dass die zur Vorlage gehörenden Dateien in einem eigenen Verzeichnis (\verb|template|) liegen.
\item Stellenweise wurden Erläuterungen als Kommentare hinzugefügt.
\end{itemize}

\item \verb|_dhbw_biblatex-config.tex|
\begin{itemize}
\item Code, der mehrere Quellenverzeichnisse unterstützt, wurde entfernt.
\item Ein zu großer Abstand nach Zitaten von Sekundärliteratur wurde korrigiert. 
\end{itemize}

\item \verb|_dhbw_erklaerung.bbx|
\begin{itemize}
\item Gemäß der Anforderung in den Zitierrichtlinien wird die Erklärung nicht ins Inhaltsverzeichnis aufgenommen und nicht mit einer Seitenzahl versehen. 
\end{itemize}

\pagebreak
\item \verb|_dhbw_praeambel.bbx|
\begin{itemize}
\item Gemäß der Anforderung in den Zitierrichtlinien werden im Literaturverzeichnis alle Autor/innen eines Werks angegeben.
\end{itemize}

\item \verb|abstract.tex|
\begin{itemize}
\item Hinweis auf \LaTeX-Spickzettel hinzugefügt.
\end{itemize}

\item \verb|deckblatt.tex|
\begin{itemize}
\item Vorname, Name, Titel der Arbeit sind nur zu Beginn einzutragen und werden dann an den entsprechenden Stellen automatisch ergänzt.
\item Hervorhebung, dass Angaben zum Unternehmen sowie den Betreuer/innen zu ergänzen sind. 
\item Wortlaut des Vertraulichkeitsvermerks wurde an die aktuelle Fassung in der Studien- und Prüfungsordnung angepasst. 
\end{itemize}

\item \verb|einleitung.tex|
\begin{itemize}
\item Ein eigenständiges Gesprächsverzeichnis als Teil des Quellenverzeichnisses ist in den Zitierrichtlinien nicht mehr vorgesehen, die entsprechenden Hinweise wurden entfernt.
\item Ein alter Hinweis auf die Darstellung von Links im Verzeichnis der Internetquellen wurde entfernt, da es ein solches eigenständiges Verzeichnis nicht mehr gibt. 
\end{itemize}

\item \verb|text_mit_zitaten.tex|
\begin{itemize}
\item Es wird nun erläutert, wie zwei Quellenangaben unmittelbar nebeneinander dargestellt werden können.
\item Erklärungen, die von mehreren Quellenverzeichnissen ausgegangen sind, wurden entfernt.
\end{itemize}

\item \verb|literatur-datenbank.bib|
\begin{itemize}
\item Gespräch wurde entfernt, da dieses nicht mehr im Quellenverzeichnis aufgeführt werden soll.
\end{itemize}

\end{itemize}
%%% Ende des eigentlichen Inhalts %%%


%%% Quellenverzeichnisse (keine Anpassung nötig) %%%
\clearpage
\literaturverzeichnis
%%% Ende Quellenverzeichnisse %%%


%%% Erklärung (keine Anpassungen nötig) %%%
% steht ganz am Ende des Dokuments
\cleardoublepage
\clearpage

\thispagestyle{empty}

{\LARGE\textsf{\textbf{Erklärung}}\bigskip}

% \typMeinerArbeit und \themaMeinerArbeit werden in deckblatt.tex definiert
Ich versichere hiermit, dass ich meine \typMeinerArbeit\ mit dem Thema: \emph{\themaMeinerArbeit} selbstständig verfasst und keine anderen als die angegebenen Quellen und Hilfsmittel benutzt habe.
Ich versichere zudem, dass die eingereichte elektronische Fassung mit der gedruckten Fassung übereinstimmt.

\vspace{3cm}

\begin{center}
\begin{tabular}{ccc}
(Ort, Datum) & \hspace{0.3\linewidth} & (Unterschrift)
\end{tabular}
\end{center}
\end{document}